%----------------------Нүүр хуудастай хамаатай зүйлс----------------------------
\pagenumbering{roman}
\makefrontpage
\maketitle

\doublespace

% Decleration
\begin{huge}
\textbf{Зохиогчийн баталгаа}
\end{huge}
\vspace{1em}
\doublespace
Миний бие \@author \ "\@title" \ сэдэвтэй судалгааны ажлыг гүйцэтгэсэн болохыг зарлаж дараах зүйлсийг баталж байна:
\begin{itemize}
\item Ажил нь бүхэлдээ эсвэл ихэнхдээ Монгол Улсын Их Сургуулийн зэрэг горилохоор дэвшүүлсэн болно.
\item Энэ ажлын аль нэг хэсгийг эсвэл бүхлээр нь ямар нэг их, дээд сургуулийн зэрэг горилохоор оруулж байгаагүй.
\item Бусдын хийсэн ажлаас хуулбарлаагүй, ашигласан бол ишлэл, зүүлт хийсэн.
\item Ажлыг би өөрөө (хамтарч) хийсэн ба миний хийсэн ажил, үзүүлсэн дэмжлэгийг дипломын ажилд тодорхой тусгасан. 
\item Ажилд тусалсан бүх эх сурвалжид талархаж байна. 
\end{itemize}
\vspace{1em}

Гарын үсэг: \underline{\hspace{5cm}}

Огноо: \underline{\hspace{3cm}}

% Гарчгийг автоматаар оруулна
\setcounter{tocdepth}{1}
\tableofcontents

% Зургийн жагсаалтыг автоматаар оруулна
\listoffigures

% Хүснэгтийн жагсаалтыг автоматаар оруулна
\listoftables

% Кодын жагсаалтыг автоматаар оруулна
\lstlistoflistings

% This puts the word "Page" right justified above everything else.
\newpage
%% \addtocontents{lof}{Зураг~\hfill Хуудас \par}
\newpage
%% \addtocontents{lot}{Хүснэгт~\hfill Хуудас \par}

\renewcommand{\cftlabel}{Зураг}

%----------------------------------------------------------------------------------------
%   УДИРТГАЛ
%----------------------------------------------------------------------------------------
\newpage
\addcontentsline{toc}{part}{УДИРТГАЛ}
\chapter*{УДИРТГАЛ}

Орчин үеийн жижиглэн худалдааны салбар нь өдөр бүр сая сая бараа бүтээгдэхүүнийг хэрэглэгчдэд хүргэж байдаг бөгөөд энэхүү үйл ажиллагааг үр дүнтэй явуулахад бараа материалын нарийвчлалтай тооллого, хяналт шалгалт чухал үүрэг гүйцэтгэдэг. Гэсэн хэдий ч уламжлалт гар аргаар хийгддэг тооллого, аудит нь цаг хугацаа их шаарддаг, хүний алдаанд өртөмтгий, үр ашиг багатай байдаг нь салбарын нийтлэг бэрхшээл болсоор иржээ.

Дэлхийн жижиглэн худалдааны зах зээл 2023 онд 28 их наяд ам.долларт хүрсэн бөгөөд бараа материалын буруу удирдлагаас үүдэлтэй алдагдал жилд 1.1 их наяд ам.долларт хүрч байна \cite{retail_stats}. Энэхүү асуудлыг шийдвэрлэхэд компьютерын хараа (computer vision) болон хиймэл оюун ухааны технологи шинэ боломжуудыг нээж байна.

\textbf{Судалгааны зорилго:} Энэхүү бакалаврын судалгааны ажлын зорилго нь дэлгүүрийн лангуун дээрх бараа бүтээгдэхүүнийг зурган дээрээс автоматаар таньж, тооллого болон аудитын шийдвэр гаргалтыг автоматжуулсан төгсгөл-хоорондын (end-to-end) систем хөгжүүлэхэд оршино.

\textbf{Судалгааны зорилтууд:}
\begin{enumerate}
    \item Жижиглэн худалдааны бараа таних YOLOv8 суурилсан объект илрүүлэлтийн загвар сургах
    \item FastAPI дээр суурилсан REST API backend сервер хөгжүүлэх
    \item Мобайл болон вэб интерфэйс хөгжүүлэх
    \item Аудитын автомат шийдвэр гаргах алгоритм боловсруулах
    \item Системийн нарийвчлал, хурдыг туршиж үнэлэх
\end{enumerate}

\textbf{Судалгааны ач холбогдол:} Хэрэгжүүлсэн систем нь жижиглэн худалдааны байгууллагуудад дараах үр өгөөжийг авчрах боломжтой:
\begin{itemize}
    \item Бараа материалын тооллогын хугацааг 80\%-иар багасгах
    \item Хүний алдааг бууруулж, нарийвчлалыг сайжруулах
    \item Бодит хугацаан дахь мэдээлэл олж авах боломж бүрдүүлэх
    \item Хөдөлмөрийн зардлыг хэмнэх
\end{itemize}

\textbf{Судалгааны арга зүй:} Судалгааг хийхдээ дараах арга зүйг баримталсан:
\begin{itemize}
    \item Онолын судалгаа: Компьютерын хараа, объект илрүүлэлт, YOLO архитектурын талаарх эрдэм шинжилгээний бүтээлүүдийг судлах
    \item Туршилтын судалгаа: Бодит дэлгүүрийн лангуунаас зураг цуглуулж, загвар сургаж, системийг туршин үнэлэх
    \item Хөгжүүлэлтийн арга зүй: Agile/Scrum арга зүйг баримтлан системийг давтамжтай хөгжүүлэх
\end{itemize}

\textbf{Бүтээлийн бүтэц:} Энэхүү судалгааны ажил нь дараах бүлгүүдээс бүрдэнэ:
\begin{itemize}
    \item \textbf{1-р бүлэг:} Оршил — судалгааны үндэслэл, зорилго, зорилтууд
    \item \textbf{2-р бүлэг:} Онолын судалгаа — объект илрүүлэлт, YOLO архитектур
    \item \textbf{3-р бүлэг:} Системийн шинжилгээ, зохиомж — архитектур, дизайн
    \item \textbf{4-р бүлэг:} Технологийн сонголт — ашигласан технологиуд
    \item \textbf{5-р бүлэг:} Хэрэгжүүлэлт — датасет, сургалт, код
    \item \textbf{6-р бүлэг:} Үр дүн, үнэлгээ — туршилтын үр дүн
    \item \textbf{7-р бүлэг:} Дүгнэлт — ерөнхий дүгнэлт, цаашдын төлөвлөгөө
\end{itemize}

%----------------------------------------------------------------------------------------
%   НЭР ТОМЪЁОНЫ ТАЙЛБАР
%----------------------------------------------------------------------------------------
\newpage
\addcontentsline{toc}{part}{НЭР ТОМЪЁОНЫ ТАЙЛБАР}
\chapter*{НЭР ТОМЪЁОНЫ ТАЙЛБАР}

\begin{longtable}{|p{4.5cm}|p{10.5cm}|}
\hline
\textbf{Нэр томъёо} & \textbf{Тайлбар} \\
\hline
\endfirsthead
\hline
\textbf{Нэр томъёо} & \textbf{Тайлбар} \\
\hline
\endhead
\hline
\endfoot

YOLO & You Only Look Once — нэг дамжуулалтаар объект илрүүлдэг гүн сургалтын архитектур \\
\hline
YOLOv8 & YOLO архитектурын хамгийн сүүлийн үеийн хувилбар (Ultralytics, 2023) \\
\hline
CNN & Convolutional Neural Network — конволюцийн мэдрэлийн сүлжээ \\
\hline
mAP & Mean Average Precision — объект илрүүлэлтийн нарийвчлалын хэмжүүр \\
\hline
IoU & Intersection over Union — хоёр хайрцагны давхцлын харьцаа \\
\hline
NMS & Non-Maximum Suppression — давхардсан илрүүлэлтийг арилгах алгоритм \\
\hline
Bounding Box & Объектыг хүрээлсэн тэгш өнцөгт хайрцаг \\
\hline
Anchor Box & Урьдчилан тодорхойлсон хайрцагны хэлбэрүүд \\
\hline
Backbone & Онцлог ялгах үндсэн сүлжээ (жнь: CSPDarknet) \\
\hline
Neck & Backbone-аас авсан онцлогуудыг нэгтгэх давхарга (жнь: PANet) \\
\hline
Head & Эцсийн таамаглал хийх давхарга \\
\hline
FPN & Feature Pyramid Network — олон масштабын онцлогуудыг нэгтгэх сүлжээ \\
\hline
Transfer Learning & Урьд сургасан загварын мэдлэгийг шинэ даалгаварт ашиглах \\
\hline
Data Augmentation & Өгөгдөл баяжуулалт — сургалтын өгөгдлийг хувиргаж олшруулах \\
\hline
Epoch & Сургалтын нэг бүрэн давталт \\
\hline
Batch Size & Нэг удаад боловсруулах зургийн тоо \\
\hline
Learning Rate & Суралцах хурд — градиент уруудалтын алхамын хэмжээ \\
\hline
Overfitting & Хэт сургалт — загвар сургалтын өгөгдөлд хэт тохирсон \\
\hline
REST API & Representational State Transfer — вэб үйлчилгээний архитектур \\
\hline
FastAPI & Python-д зориулсан өндөр хурдтай вэб фреймворк \\
\hline
MongoDB & NoSQL баримт бичигт суурилсан өгөгдлийн сан \\
\hline
Docker & Контейнержуулалтын платформ \\
\hline
Odoo & Нээлттэй эхийн ERP систем \\
\hline
Audit & Тооллого, шалгалт — бараа материалын бодит байдлыг шалгах үйл явц \\
\hline
Planogram & Лангууны байршлын зураглал \\
\hline
SKU & Stock Keeping Unit — барааны нэгжийн код \\
\hline
Real-time & Бодит хугацаанд — хоцролтгүй шууд боловсруулалт \\
\hline
End-to-end & Төгсгөл-хоорондын — эхнээс төгсгөл хүртэл бүрэн систем \\
\hline
Precision & Нарийвчлал — зөв эерэг / нийт эерэг таамаглал \\
\hline
Recall & Бүрэн байдал — зөв эерэг / нийт бодит эерэг \\
\hline
\end{longtable}

\doublespace
\pagenumbering{arabic}
